% SOURCE_AUTHORITY: sga5_fr_workpass.tex, lines 14205--14240 (LF-normalized unit).
% SOURCE_SHA256: 791F4EFFC5E02832D5D77ED03518C8156D6F07E4C8238B03545DB93D883FBB28
\begin{statement}{Exemplo 7.3}
Sejam $X$ um pré-esquema definido sobre o corpo finito $\mathbb F_p$, onde $p$ é um número primo, $\fr_X:X\to X$ o endomorfismo de Frobenius (cf. XIV), e $F$ um feixe de $\Lambda$-módulos sobre $X$.

Se $Y,S$ são esquemas sobre $\mathbb F_p$, e se $Y$ é um esquema sobre $S$, denotamos por
\[
\Fr_{Y/S}:Y\longrightarrow Y^{(p/S)}
\]
o morfismo de Frobenius relativo (cf. XIV 2).

Obtém-se um morfismo de feixes sobre $X_{\et}$
\[
(\fr_X)_*(F)\longrightarrow F
\]
colocando, para todo $X$-pré-esquema étale $U$,
\[
(\fr_X)_*(F)(U)=F(\fr_X^{-1}(U))=F(U^{(p/X)})
 \xrightarrow{\,F(\Fr_{U/X})\,}F(U).
\]
Este morfismo é, de fato, um isomorfismo. Seu inverso
\[
F(\Fr_{/X})^{-1}:F\longrightarrow (\fr_X)_*(F)
\]
define por adjunção um morfismo
\[
\Fr_{F/X}^*:(\fr_X)^*(F)\longrightarrow F
\]
que é um isomorfismo.

Se, em particular, $X$ é suave de dimensão $1$ e $F$ construtível, o par
\[
(\Fr_{F/X}^*,\fr_X):(X,F)\longrightarrow (X,F)
\]
fica, portanto, abrangido pela proposição 7.2.
\end{statement}

\noindent[1] Wecken, F., Fixpunktklassen I, Math. Ann. 117 (1941), p. 659--671; Fixpunktklassen II, III, Math. Ann. 118 (1942-1943), p. 216-234 e p. 544-577.
